\documentclass[12pt]{article}

\usepackage[margin=1in]{geometry}
\usepackage[T1]{fontenc}
\usepackage{lmodern}
\usepackage{microtype}
\usepackage[hidelinks]{hyperref}

\title{Food Displays Play a Key Role in Enhancing Dietary Diversity Among Farming Households in Chikwawa District}
\author{Aaron Posiano \and Kenani Kalagho\\
\small Shire Valley Agricultural Development Division, P/Bag 1, Ngabu}
\date{}

\begin{document}
\maketitle

\noindent\textbf{Background:}
Dietary diversity is a key indicator of diet quality and overall nutritional status. Strategic food displays can make healthy choices more visible and appealing, thereby encouraging experimentation and adoption of balanced diets. In farming households, where food choices are influenced by availability and cultural norms, visual cues such as food displays can be a powerful behavior change tool.

\medskip
\noindent\textbf{Objective:}
The study aimed to identify key factors that influence dietary diversity among farming households in 7 EPAs of Chikwawa District.

\medskip
\noindent\textbf{Methodology:}
Cross section data was collected from 104 farming households using a structured questionnaire in a routine surveillance exercise in the 7 EPAs of Chikwawa District. Descriptive (frequencies) and hypothesis testing (multiple regression) analytical techniques were undertaken using Microsoft Excel and R respectively at 5\% significance.

\medskip
\noindent\textbf{Results:}
Of the respondents, 54.8\% achieved medium dietary diversity (MDD), 25.0\% achieved high dietary diversity (HDD), and 20.2\% had low dietary diversity (LDD). In five of the seven EPAs, at least half of the respondents attained MDD, while two EPAs showed a mix of HDD and LDD. Multiple regression analysis identified participation in nutrition training ($p = 0.04414$) and exposure to food displays ($p = 0.00864$) as significant predictors of dietary diversity. The model's $F$-statistic (3.31) exceeded the tabulated value (1.75), with a $p$-value of 0.0002171, confirming the significant influence of these interventions on dietary diversity scores.

\medskip
\noindent\textbf{Conclusion:}
Participation in nutrition training and exposure to food displays significantly enhance dietary diversity among farming households. These interventions promote the consumption of diverse, locally available foods, contributing to improved nutritional outcomes.

\medskip
\noindent\textbf{Recommendation:}
Nutrition education efforts should incorporate displays of actual and locally sourced food models. Scaling up food displays across EPAs and building the capacity of frontline staff to deliver engaging demonstrations can promote long-term dietary improvements within farming communities.

\medskip
\noindent\textbf{Keywords:} Dietary Diversity, Food Displays, Nutrition Education, farming households

\medskip
\noindent\textbf{Lead Author Contacts:}
\href{mailto:aaronposiano@gmail.com}{aaronposiano@gmail.com},
+265(0)881124453/993524636

\end{document}